% [Sumaiya] Related Work (rubric: 2 marks).
% Top band: AT LEAST 5 scholarly sources, COMPARED and explicitly POSITIONED against this work.
%   3-4 sources = Good; 1-2 listed = Satisfactory. Not a list: say what each did and how we differ.
% Sources come from everyone (T1, two each in refs.bib by CP2). Suggested groups:
%   ETA / delivery-time prediction; the Meituan dataset paper; quantile regression / promise setting;
%   leakage in ML (Kapoor & Narayanan 2023, cited in Lec 6).
\section{Related Work}

Predicting delivery time from data has been studied mostly with rich platform data. Zhu
\textit{et al.}~\cite{zhu2020fulfillment} estimate the full order-to-delivery time for Ele.me and deploy
the model there; theirs is the closest task to ours, but it draws on proprietary platform data that are
not public. Hildebrandt and Ulmer~\cite{hildebrandt2022supervised} study arrival-time estimation in
restaurant meal delivery, the same prediction problem: they map the current state of the delivery fleet
to arrival times with supervised learning, and match a near-optimal online simulation at a fraction of its
cost. They assume the dispatcher's live state is available; we have only what is known about the order
itself. Wang
\textit{et al.}~\cite{wang2018travel} learn travel times for DiDi ride-hailing at industrial scale; food
delivery adds a stage before the ride (dispatch, reaching the restaurant and waiting for the food), which
our stage analysis separates from the road time. We differ from all three in working with a public log,
released for the INFORMS TSL challenge~\cite{zhao2024meituan}, with the much smaller set of order-time
features such a log provides, and in measuring what the missing after-order information is worth
(Section~\ref{sec:leakage}).

A second line of work uses predictions to act. Liu \textit{et al.}~\cite{liu2021ontime} feed predicted
travel times into order assignment to improve on-time delivery. Ulmer \textit{et al.}~\cite{ulmer2021restaurant}
route couriers when restaurant ready times are random. Both act on dispatch, deciding who delivers and
when; we act on the promise shown to the customer, and compare an 80th-percentile promise with the
platform's own at the same average length. Our audit adds that 13.8\% of the platform's ready-time
estimates are mostly system defaults, not forecasts.

Finally, Kapoor and Narayanan~\cite{kapoor2023leakage} find leakage in 294 papers across 17 fields; one
of its forms is information unavailable at prediction time entering a model. Delivery logs invite it,
since dispatch, pickup and arrival times sit next to the order, so we keep after-order times out of every
model and measure their worth separately.

Deployed systems (Table~\ref{tab:benchmark} in the appendix) use the same kinds of real-time inputs as
our extension (Section~\ref{sec:extension}): Ele.me's model~\cite{zhu2020fulfillment,alibaba2020fulfillment}
cut MAE by 9.8\% online with them, and Uber's DeepETA~\cite{uber2022deepeta,hu2022deepreta} uses an
asymmetric loss so that being late costs more, the trade-off a padded promise makes.
